Euler's method is a straightforward numerical technique for solving differential equations by constructing a continuous set of approximations using tangent lines computed at discrete points.
Given a differential equation with an initial condition:

\[
\frac{dy}{dt} = f(t, y), \quad y(t_0) = y_0,
\]

we start at the point \((t_0, y_0)\) and use the slope to form the tangent line and compute subsequent approximations:

\[
y = y_0 + f(t_0, y_0)(t - t_0)
\quad \rightarrow \quad
y_{n+1} = y_n + f(t_n, y_n) \cdot \Delta t,
\]

where \(\Delta t = t_{n+1} - t_n\) is the fixed step size.

The first equation gives the tangent line approximation from the initial point.
Evaluating this at \(t = t_1\) provides the next point \((t_1, y_1)\).
The iterative process continues using the second equation for all subsequent steps.
As the number of steps increases and \(\Delta t\) decreases, this piecewise linear approximation converges to the true solution \(\phi(t)\).